% Continuous sampling models, input from Chapter 5.
\section{Bertrand's three experiments}

A chord in a circle is longer than the side of an inscribed equilateral triangle when its midpoint lies inside the concentric disk of half the radius. Three different experiments give three probabilities. Joseph Bertrand set the question in \emph{Calcul des probabilit\'es} (1889), page~4, and refused to choose. Take the radius to be $1$. The side of the inscribed equilateral triangle subtends the angle $2\pi/3$ at the centre.

\begin{figure}[ht]
\centering
\fig{ch01-bertrand}
\caption{Three experiments, one circle of radius $1$. Left: one endpoint is fixed and the shorter angle $\alpha$ to the other endpoint exceeds $2\pi/3$ on a third of the circle. Middle: the distance $D$ from the centre to the chord is drawn uniformly from $[0,1]$, and the long chords are those with $D<1/2$. Right: the midpoint is uniform in the disk, and the long chords are those whose midpoint falls in the inner disk, of a quarter of the area.}
\label{fig:bertrand}
\end{figure}

\begin{theorem}[three models for a chord]\label{thm:bertrand}
\begin{enumerate}[label=(\alph*)]
\item If the two endpoints are independent and uniform on the circle, the probability that the chord is long is $\frac13$.
\item If a radius is chosen, and the distance from the centre to the chord is then uniform on $[0,1]$ and perpendicular to that radius, the probability is $\frac12$.
\item If the midpoint is uniform in the disk, the probability is $\frac14$.
\end{enumerate}
\end{theorem}
\begin{proof}
(a) Rotational symmetry lets us fix one endpoint. The other endpoint is uniform on the circle. Let $\alpha\in(0,\pi]$ be the shorter angle between the endpoints. For $\theta\in(0,\pi]$, the event $\alpha\le\theta$ means that the free endpoint lies in one of the two arcs of length $\theta$ adjacent to the fixed endpoint, a fraction $\theta/\pi$ of the circle. So $P(\alpha\le\theta)=\theta/\pi$, and $\alpha$ has density $1/\pi$ on $(0,\pi]$. The chord is long when $\alpha>2\pi/3$, which has probability $(\pi-2\pi/3)/\pi=1/3$.

(b) In this model the distance $D$ is uniform on $[0,1]$ by the description of the experiment, not by a theorem about endpoints. The chord is long when $D<1/2$, which has probability $1/2$.

(c) The set of midpoints that give a long chord is the open disk of radius $1/2$. Its area is $\pi/4$. The whole disk has area $\pi$. A uniform midpoint falls in the inner disk with probability $1/4$. The area of a disk is the integral computed in Volume~1, Chapter~8: the quarter circle $y=\sqrt{1-x^2}$ on $[0,1]$ has area $\pi/4$.
\end{proof}

The three numbers disagree because the three sets of outcomes disagree. Model (a) spreads endpoints. Model (b) spreads distances along a radius. Model (c) spreads area. Bertrand's sentence is the instruction this book follows. Write the experiment, then compute.

\begin{examplex}[the length, under one model]
Keep the radius equal to $1$, and use model (b): the distance $D$ from the centre to the chord is uniform on $[0,1]$. The chord length is $L=2\sqrt{1-D^2}$. It falls from $2$ to $0$ as $D$ runs from $0$ to $1$. For $\ell\in(0,2)$,
\[
 P(L\le\ell)=P\Bigl(D\ge\sqrt{1-(\ell/2)^2}\Bigr)=1-\sqrt{1-(\ell/2)^2}.
\]
Differentiate. The chain rule gives the density
\[
 f_L(\ell)=\frac{\ell}{4\sqrt{1-(\ell/2)^2}}\qquad\text{on }(0,2).
\]
The long chords are those longer than the side $\sqrt3$ of the inscribed equilateral triangle. The formula yields $P(L>\sqrt3)=1/2$, which is clause (b) of the theorem again: $\sqrt{1-3/4}=1/2$.

The density blows up as $\ell\uparrow2$ and tends to $0$ as $\ell\downarrow0$. Uniform distance spends as much probability on $D<1/2$ as on $D>1/2$, and $D<1/2$ is exactly the long chords. Area measure does not. The disk of radius $1/2$ has a quarter of the area, which is why model (c) gives probability $1/4$ for the same geometric event. The map from distance to length is not the identity, so a uniform law for $D$ is not a uniform law for $L$, and neither of those is the area law for the midpoint.
\end{examplex}

\section{Buffon's needle}

Parallel lines are ruled on the plane at distance $d>0$. A needle of length $\ell$, with $0<\ell\le d$, is thrown. Buffon described this experiment in the \emph{Essai d'arithm\'etique morale}, in the supplement to the \emph{Histoire naturelle}, in 1777. The model we adopt is the following. Let $X$ be the distance from the needle's centre to the nearest line, uniform on $[0,d/2]$. Let $\Theta$ be the angle between the needle and the direction of the lines, taken in $[0,\pi]$ by turning the needle end for end, and take $\Theta$ uniform on $[0,\pi]$ and independent of $X$. The needle meets a line when the vertical half-projection reaches the nearest line, that is when $X\le(\ell/2)\sin\Theta$. The restriction $\ell\le d$ keeps the needle from meeting two lines at once, so this inequality is the whole event.

\begin{figure}[ht]
\centering
\fig{ch01-buffon}
\caption{Two rulings at distance $d$. The needle has length $\ell\le d$. It meets a ruling when the distance $X$ from its centre to the nearest ruling is at most $(\ell/2)\sin\Theta$.}
\label{fig:buffon}
\end{figure}

\begin{theorem}[Buffon's needle]\label{thm:buffon}
Under the model just stated, the probability that the needle meets a line is $2\ell/(\pi d)$.
\end{theorem}
\begin{proof}
Independence and uniformity put density $2/d$ on $X\in[0,d/2]$ and density $1/\pi$ on $\Theta\in[0,\pi]$. The probability is the integral of the product density over the region $0\le x\le(\ell/2)\sin\theta$. For each fixed $\theta$ the inner integral in $x$ is an ordinary one-variable integral, and $(\ell/2)\sin\theta\le d/2$ because $\ell\le d$ and $\sin\theta\le1$. So
\begin{align*}
P
&=\int_0^\pi\frac1\pi\Biggl(\int_0^{(\ell/2)\sin\theta}\frac2d\,\mathrm dx\Biggr)\mathrm d\theta
=\frac{2}{d\pi}\int_0^\pi\frac\ell2\sin\theta\,\mathrm d\theta
=\frac{\ell}{d\pi}\bigl[-\cos\theta\bigr]_0^\pi
=\frac{2\ell}{\pi d}.
\end{align*}
The iterated integral of a continuous function on a rectangle, cut down to a region described by $\sin$, is justified by writing it as a single integral of a continuous function of $\theta$, which Volume~1, Chapter~8, allows. The repeated-integral statement for a continuous integrand on a rectangle is Proposition~\ref{prop:fubini} above.
\end{proof}

When $\ell=d$, the probability is $2/\pi\approx0.6366$. A needle as long as the gap hits a line more often than not, and not twice as often as a needle of half the length would suggest if one forgot the angle. The factor $\sin\Theta$ is the whole geometry.

The same integral is the first geometric probability in the series. Volume~9 returns to it for a curve rather than a needle. The proof there may use more geometry. It does not replace the computation above.
